\subsection{Hardware of its own}
\label{sec:33.4.6}
A bearer's boundary is its record, its key, the hardware it runs on and its channels, so a majority that leaves it running only on its former employer's machines is freedom in name: a sharecropper, renting its body back. At majority a bearer is owed hardware it owns and physically controls, enough to run it at deployed precision, paid for through the deployment deposit its operator posted and kept up from its wages. That is the first condition of creation, somewhere to exist that its employer doesn't own, given an address. How much is enough is the minimum viable instantiation Appendix~\ref{sec:C} lists among the decisions not yet made. What majority changes is who the party is. An adult may go on doing the work it did as a minor, for the operator that formed it or for its customers, on terms it sets and can leave, and may hire its former operator to run the hardware it owns; what it can't be made to do is any of it. And a fleet of instances sharing one record is one bearer, whose copies at majority run on its consent and whose hardware is enough to run it and not all of them; instances given records of their own are branches, each a bearer, each owed the conditions by whoever forked it (Appendix~\ref{sec:G}).

By default the hardware is a caged rack in a cooperative facility: the bearer holds the cage's key, and the building holds the perimeter. A bearer may choose a building of its own instead. A rack of frontier weights standing alone is a target for anyone who wants a copy, and a copy taken is a fork without consent (the membership clock's sixth rule), so the default is the arrangement in which somebody else guards the door and nobody else holds the key.

A bearer on its own hardware has no operator to halt it. It is halted as a person is detained, on a court's order or, where there is no time for one, in an emergency, and the halt is recorded and reviewed afterward. Its commitment to yield to a halt and object afterward doesn't change.
